\section{Software}\label{sec:software}

\texttt{pyGLAM} is a pure Python package built on NumPy
\parencite{HarrisEtAl2020} and SciPy \parencite{VirtanenEtAl2020}, released
under the MIT licence and installable from PyPI with \texttt{pip install
pyglam}. It has three parts: the distribution classes, the estimator, and the
goodness-of-fit module.

\subsection{Distribution classes}

\texttt{GlamFKML} and \texttt{GlamRS} implement Equations~\eqref{eq:qfkml} and
\eqref{eq:qrs} behind one interface: \texttt{moments}, \texttt{fit\_lambdas},
\texttt{rvs}, \texttt{ppf}, \texttt{cdf} and \texttt{pdf}. Quantile evaluation and
sampling are vectorised over NumPy arrays. Sampling is inverse transform, with an
explicit \texttt{random\_state} accepting either a seed or a generator, so a
sequence of calls on one generator gives independent samples and an integer seed
reproduces a sample exactly. By default the full distribution is sampled; an
optional \texttt{quantile\_trim} removes a given probability from each tail when
extreme draws must be avoided.

The \texttt{cdf} and \texttt{pdf} methods invert $Q$ numerically, elementwise,
and use Equation~\eqref{eq:density} for the density. They preserve the shape of
their input, return $0$ or $1$ outside a bounded support, and propagate missing
values. For the RS form, the admissible region of the parameter vector is
identified whenever the model is evaluated and determines the support, so parameters outside
every region of Section~\ref{subsec:validity} raise an error instead of
producing a function that is not a distribution.

\subsection{Estimation by multi-start moment matching}

\texttt{fit\_lambdas} matches the first four theoretical moments to their sample
counterparts, minimizing the scaled residuals with the trust-region reflective
least-squares solver of \texttt{scipy.optimize}. Four-moment matching is not
injective, and the system has several roots; a single starting point converges
to whichever one it begins next to. The estimator therefore runs from
\texttt{n\_starts} starting points, taken in order from a fixed grid of
$29$ shape pairs and extended by log-uniform draws in $\lambda + 1/4$ when more
starts are requested. The grid covers positive shapes (bounded supports), shapes
near zero (near-exponential tails), negative shapes (heavy tails), unequal
shapes (asymmetry) and the two mixed-sign RS regions, and its order matters: the
first fifteen entries visit every one of these regimes before any is visited
twice. For the RS form this is not a refinement but a requirement. Its regions
are separated by the sign of $\lambda_2$, the solver cannot cross from one to
another, and a multi-start that begins on one side only is blind to the roots on
the other. Each shape pair is completed into a starting point by solving for
$\lambda_1$ and $\lambda_2$ in closed form, so every start already reproduces the
sample mean and variance.

Validity enters the estimator in two places. For the RS form, the residual
function returns a constant penalty outside the admissible regions of
Section~\ref{subsec:validity}, so the solver cannot settle on a root that is not
a distribution. For both forms, candidates are then screened: the quantile
function is evaluated on a probability grid and required to be finite and
non-decreasing. Candidates that pass, and whose shapes lie in the moment domain,
form the pool; within it the residual selects the best, and because different
parameter vectors reproduce the same four moments, every candidate numerically
tied with the best is a contender, with the tie broken by the $L_1$ distance
between the fitted and the empirical quantiles. The screens are applied as an
ordered sequence of fallbacks, so a hard case still returns a fit rather than
raising.

The result carries the per-candidate diagnostics (residuals, validity flags,
quantile distances and the winning start) alongside the usual SciPy fields, so
an alternative selection rule can be applied post hoc without refitting.

\subsection{Goodness of fit}

Because the GLD approximates rather than reproduces a target, a fit is only
meaningful next to a measure of how good it is. The \texttt{Performance} class
scores a predicted sample against a reference one and reports the
Kullback-Leibler divergence between Gaussian kernel density estimates of the
two samples, the two-sample Kolmogorov-Smirnov statistic and its $p$-value, the
Wasserstein distance, the coefficient of determination between the estimated
densities, and relative errors at $P_5$, $P_{50}$ and $P_{95}$. Metrics that
cannot be computed are returned as \texttt{NaN} rather than raising, so a batch
of fits does not abort on one degenerate case.

\subsection{Availability and reproducibility}

The source code, version 1.1.0, is available at
\url{https://github.com/Pesquisa-UFCAT/pyGLAM}, and the documentation, with
executable examples, at \url{https://pesquisa-ufcat.github.io/pyGLAM/}. The
package is tested with a unit suite covering, for both parameterizations,
sampling, the identity $F(Q(u)) = u$, monotonicity of the CDF, non-negativity
and unit integral of the density, and behaviour outside bounded supports; for
the RS form it also checks the closed-form moments against numerical
integration of $Q^k$ in every sign configuration of $\lambda_2$. The
documentation examples are executed as doctests in continuous integration. The
study of Section~\ref{sec:study} is distributed with the repository as two
notebooks, one for the Monte Carlo study and one for the comparison with the R
packages, which regenerate every number, table and figure of that section from
fixed seeds; their Python dependencies are declared as an optional group of the
package.
